\chapter{Reduction from SBI to the bandit with expert advice (Section 3)}
\label{chap:reduction}

Throughout, $0 < \eps \le 0.1$, and $u^\star$ and $j^\star$ denote the special batch and the
\emph{special expert} of the strategy $S_\theta$: $j^\star = \theta$ (expert $0$ for $S_0$).
$N_u(T)$ is the number of rounds $t < T$ in which an arm of batch $u$ is pulled.

The lemmas of this chapter are formalized on the general arm and expert types of
Chapter~\ref{chap:main} (\texttt{strategyOn}, Definition~\ref{def:embedding}): an arm $a$ is in
the batch of its reduction $\sigma(a)$, and the special expert is $\tau'(\theta)$.

\begin{lemma}[Conditional law of the losses]
  \label{lem:loss_law_cond}
  \lean{Chase2026Tight.hasCondDistrib_feedback_strategyOn, Chase2026Tight.lossKernelOn}
  \leanok
  \uses{def:strategy, lem:feedback_ignores_action_cond}
  In a run of a learner against $S_\theta$, the loss vector of round $t$ has conditional law
  $\texttt{lossLaw ε θ}(e_t)$ given the advice $e_t$ and the arm $I_t$ of round $t$ (and also given
  the history of the previous rounds, the advice and the arm, by definition of a run).
\end{lemma}

\begin{proof}
  \leanok
  The feedback kernel of $S_\theta$ is the loss kernel applied to the advice of the round, a
  function of the advice and the arm; Lemma~\ref{lem:feedback_ignores_action_cond}.
\end{proof}

\begin{lemma}[Expected losses given the advice]
  \label{lem:mean_loss}
  \lean{Chase2026Tight.meanLoss, Chase2026Tight.meanLoss_zero,
    Chase2026Tight.meanLoss_ofBatch_of_ne, Chase2026Tight.meanLoss_mk_ofBatch,
    Chase2026Tight.meanLoss_sub_meanLoss_adviceOf_ge}
  \leanok
  \uses{def:strategy}
  Let $0 \le \eps \le 1/2$ and let the advice $e$ be given by advice bits. Under
  $\texttt{lossLaw ε θ}(e)$, the expected loss of arm $0$ is $1/2 - \eps/2$, of an arm of a batch
  other than $u^\star$ it is $1/2$, and of the arm $(u, b)$ of the special batch of $S_{(u, v)}$ it
  is $1/2 - \eps$ if $(u, v)$ advises $(u, b)$ and $1/2 + \eps$ otherwise. Hence the expected loss
  of every arm $a$ exceeds the expected loss of the arm $e(j^\star)$ advised by the special expert
  by at least $\eps/2$ if $a \notin u^\star$, and by at least $0$ otherwise.
\end{lemma}

\begin{proof}
  \leanok
  The loss of arm $0$ is a $\mathrm{Ber}(1/2 - \eps/2)$ bit; the correct arm of a batch other than
  the special batch is a uniform bit; the arm advised by $(u, v)$ is incorrect with probability
  $1/2 - \eps$ under $S_{(u,v)}$. The advice of $j^\star$ is arm $0$ for $S_0$ and an arm of batch
  $u$ for $S_{(u, v)}$; compare the five cases.
\end{proof}

\begin{lemma}[Per-round regret]
  \label{lem:per_round_regret}
  \lean{Chase2026Tight.mul_measureReal_batch_ne_le_integral, Chase2026Tight.strategyOn_obsIn}
  \leanok
  \uses{lem:loss_law_cond, lem:mean_loss}
  In a run against $S_\theta$, for every round $t$,
  $\E[\ell_t(I_t) - \ell_t(e_t(j^\star))] \ge \frac{\eps}{2} \Prob(I_t \notin u^\star)$.
\end{lemma}

\begin{proof}
  \leanok
  Condition on the advice and the arm (Lemma~\ref{lem:loss_law_cond}): the expectation is that of
  the difference of the expected losses of $I_t$ and of $e_t(j^\star)$ given the advice. The advice
  is almost surely given by advice bits (it is drawn by \texttt{adviceOf} from uniform bits), and
  Lemma~\ref{lem:mean_loss} bounds this difference from below by $\frac{\eps}{2}
  \mathbf{1}\{I_t \notin u^\star\}$.
\end{proof}

\begin{lemma}[Regret and pulls of the special batch]
  \label{lem:regret_ge_pulls}
  \lean{Chase2026Tight.mul_sub_integral_pullCount_le_expertPseudoRegret,
    Chase2026Tight.integral_pullCount_history}
  \leanok
  \uses{lem:per_round_regret, def:pseudo_regret}
  In a run against $S_\theta$, $R_T \ge \frac{\eps}{2} \E[T - N_{u^\star}(T)]$.
\end{lemma}

\begin{proof}
  \leanok
  The pseudo-regret is at least the expected regret against $j^\star$ (the losses are in
  $\{0, 1\}$, hence integrable); sum Lemma~\ref{lem:per_round_regret} over $t < T$, using
  $\E[N_{u^\star}(T)] = \sum_{t < T} \Prob(I_t \in u^\star)$.
\end{proof}

\begin{definition}[The SBI algorithm of a learner]
  \label{def:sbi_of_learner}
  \lean{Chase2026Tight.sbiOfLearner, Chase2026Tight.mostPulledBatch}
  \leanok
  \uses{def:sbi}
  Given a learner $A$ and $T^\star \in \N$, the SBI algorithm $A'$ runs $A$ (with bandit feedback)
  for $T^\star$ rounds, then outputs a most pulled batch, an element of
  $\argmax_{u} N_u(T^\star)$ (a fixed choice among the maximizers).
\end{definition}

\begin{lemma}[Runs of the SBI algorithm of a learner]
  \label{lem:sbi_of_learner_run}
  \lean{Chase2026Tight.stoppingTime_sbiOfLearner, Chase2026Tight.isRun_sbiOfLearner,
    Chase2026Tight.mostPulledBatch_eq}
  \leanok
  \uses{def:sbi_of_learner}
  The SBI algorithm $A'$ of Definition~\ref{def:sbi_of_learner} stops after $T^\star$ rounds.
  Every run of the learner $A$, with output the most pulled batch in the first $T^\star$ rounds, is
  a run of $A'$. A batch pulled in more than half of the $T^\star$ rounds is the most pulled batch.
\end{lemma}

\begin{proof}
  \leanok
  The stopping rule is ``the history has length $T^\star$'' and the output rule is deterministic.
  Two distinct batches are pulled in disjoint sets of rounds, so if $N_u(T^\star) > T^\star/2$
  every other batch has fewer pulls than $u$.
\end{proof}

\begin{lemma}[Reduction, Lemma 3.1]
  \label{lem:reduction}
  \lean{Chase2026Tight.exists_isGood_expectedRounds_le,
    Chase2026Tight.exists_isGoodOn_lintegral_stoppingTime_le}
  \leanok
  \uses{def:is_good, def:pseudo_regret, def:sbi_of_learner, lem:sbi_of_learner_run,
    lem:regret_ge_pulls, lem:exists_run_universe}
  Let $A$ be a learner whose pseudo-regret under every strategy of the pool is at most $r(T)$ after
  $T$ rounds, for all $T$, and let $T^\star \ge 1$ with $T^\star \ge 1000\, r(T^\star)/\eps$. Then
  there is a good SBI algorithm $A'$ with $T(A', S) \le T^\star$ for every strategy $S$ of the pool.
\end{lemma}

\begin{proof}
  \leanok
  Take $A'$ of Definition~\ref{def:sbi_of_learner}: it stops at $T^\star$, and its runs are runs of
  $A$ (Lemma~\ref{lem:sbi_of_learner_run}). The law of its output under $S_\theta$ can be computed
  on any run; take a run of $A$ on a probability space in the universe of the regret hypothesis
  (Lemma~\ref{lem:exists_run_universe}). By Markov's inequality and
  Lemma~\ref{lem:regret_ge_pulls},
  $\Prob(N_{u^\star}(T^\star) \le \frac12 T^\star)
  \le \frac{2\, \E[T^\star - N_{u^\star}(T^\star)]}{T^\star}
  \le \frac{4\, r(T^\star)}{\eps T^\star} \le 0.004$, and $N_{u^\star}(T^\star) > \frac{T^\star}{2}$
  makes $u^\star$ the most pulled batch (Lemma~\ref{lem:sbi_of_learner_run}). The general form,
  on general arm and expert types, is proved in the same way; the paper's form is the case of the
  identity maps.
\end{proof}

\begin{remark}
  \label{rem:reduction_t_pos}
  The paper does not state $T^\star \ge 1$, but its proof divides by $T^\star$, and the statement
  is false for $T^\star = 0$: with $r(T) = T$ the hypotheses hold, and no good algorithm stops after
  $0$ rounds (when $k \ge 1$, the law of the output does not depend on the strategy).
\end{remark}
